% c11.fig.koordinaten — one Cartesian component of the change of coordinates.
%
% Geometry from the formulas of DERIVATION.md, Result 1.1, for one component:
%     X = (m_e x_e + m_p x_p)/M,      x = x_e - x_p,
%     x_e = X + (m_p/M) x,            x_p = X - (m_e/M) x,
%     forward matrix A = ((m_e/M, m_p/M), (1, -1)),  det A = -1,  |det A| = 1.
% Drawn with the STATED, NOT PHYSICAL ratio m_p : m_e = 3 : 1 (m_e/M = 1/4, m_p/M = 3/4):
%     lines X = k:   x_p = (4k - x_e)/3      (slope -1/3)
%     lines x = k:   x_p = x_e - k           (slope +1)
% The cell 1 <= X <= 2, 0 <= x <= 1 has the corners (x_e, x_p) = (X + 3x/4, X - x/4):
%     (1,1), (2,2), (2.75,1.75), (1.75,0.75);
% its area is |(1)(-1/4) - (1)(3/4)| = 1, the area of the unit square 0 <= x_e, x_p <= 1.
% Greyscale: the square is hatched one way, the cell the other way.
\begin{tikzpicture}[x=1.55cm, y=1.55cm, >={Stealth[length=5pt,width=4pt]}]
  \begin{scope}
    \clip (-0.75,-0.75) rectangle (3.55,2.75);
    % lines of constant X (k = -1..4) and of constant x (k = -3..4)
    \foreach \k in {-1,0,1,2,3,4}
      \draw[tpc grid, draw=tpc-plate-muted] (-1,{(4*\k+1)/3}) -- (4,{(4*\k-4)/3});
    \foreach \k in {-3,-2,-1,0,1,2,3,4}
      \draw[tpc grid, draw=tpc-plate-muted, dash pattern=on 2.5pt off 1.5pt] (-1,{-1-\k}) -- (4,{4-\k});
    % unit square of the old coordinates and unit cell of the new ones
    \draw[tpc-plate-ink, line width=\tpcstrokeregular, pattern=north west lines,
          pattern color=tpc-plate-rule] (0,0) rectangle (1,1);
    \draw[tpc-plate-ink, line width=\tpcstrokeregular, pattern=north east lines,
          pattern color=tpc-plate-rule] (1,1) -- (2,2) -- (2.75,1.75) -- (1.75,0.75) -- cycle;
  \end{scope}
  % old axes
  \draw[tpc axis, ->] (-0.75,0) -- (3.75,0) node[right, tpc label] {$x_{\mathrm{e}}$};
  \draw[tpc axis, ->] (0,-0.75) -- (0,2.95) node[above, tpc label] {$x_{\mathrm{p}}$};
  \foreach \t in {1,2,3} \draw[tpc axis] (\t,0) -- (\t,-0.06) node[below, tpc label] {$\t$};
  \foreach \t in {1,2}   \draw[tpc axis] (0,\t) -- (-0.06,\t) node[left, tpc label] {$\t$};
  % the new basis vectors at the corner (1,1): along X (x fixed) and along x (X fixed)
  \draw[draw=tpc-plate-ink, line width=\tpcstrokestate, ->] (1,1) -- (2,2);
  \draw[draw=tpc-plate-ink, line width=\tpcstrokestate, ->] (1,1) -- (1.75,0.75);
  \node[tpc label, fill=tpc-plate, inner sep=1pt, anchor=south east] at (1.52,1.60) {$\Delta X = 1$};
  \node[tpc label, fill=tpc-plate, inner sep=1pt, anchor=north] at (1.30,0.80) {$\Delta x = 1$};
  \node[tpc label, fill=tpc-plate, inner sep=1.5pt] at (0.5,0.5) {\tpcifenglish{area}{Fläche} $1$};
  \node[tpc label, fill=tpc-plate, inner sep=1.5pt] at (1.9,1.38) {\tpcifenglish{area}{Fläche} $1$};
  % labels of the two families of lines, where they leave the window (x_e = 3.55)
  \node[tpc label, anchor=west] at (3.6,{(8-3.55)/3}) {$X = 2$};
  \node[tpc label, anchor=west] at (3.6,2.55) {$x = 1$};
  \node[tpc label, anchor=north west, align=left, text=tpc-plate-muted] at (3.6,0.95)
    {\tpcifenglish{drawn for}{gezeichnet für}\\ $m_{\mathrm{p}} : m_{\mathrm{e}} = 3 : 1$};
\end{tikzpicture}
